\studytypesubsection{LLMs as Judges}
\label{sec:llms-as-judges}

\studytypeparagraph{Description}

LLMs can rate properties of software artifacts (e.g. assess code readability, adherence to coding standards, or comment quality) or sort multiple solutions by some attribute. These kinds of judgments are distinct from annotating unstructured text (see Section \annotators).

\studytypeparagraph{Example(s)}

\citeauthor{DBLP:conf/re/LubosFTGMEL24}~\cite{DBLP:conf/re/LubosFTGMEL24} leveraged \href{https://www.llama.com/llama2/}{Llama-2} to evaluate the quality of software requirements statements. 
They prompted the LLM with the text below, where the words in braces reflect the study parameters:

\judgesexample

\citeauthor{DBLP:conf/re/LubosFTGMEL24} evaluated the LLM's output against expert human judges and found moderate agreement for simple requirements and poor agreement for more complex requirements. In contrast, \citeauthor{wang2025multimodalrequirementsdatabasedacceptance}~\cite{wang2025multimodalrequirementsdatabasedacceptance} used an LLM to generate acceptance criteria for user stories, and provided a rubric to an LLM to judge the generated acceptance criteria on interpretable scales (0 to 4). 
 

\studytypeparagraph{Advantages}

Like with annotation, LLMs are often faster and cheaper than human judges, \added{so they support larger datasets for judgment studies}. However, the main constraint that varies by model and budget is the input context size, (the number of tokens one can pass into a model). For example, the upper bound of the context that can be passed to OpenAI's \textsf{o1-mini} model is \href{https://help.openai.com/en/articles/9855712-openai-o1-models-faq-chatgpt-enterprise-and-edu}{32k tokens}. Some people argue that LLMs could be more reliable or less biased than human annotators, but no one has provided any compelling evidence to support this claim.    


\studytypeparagraph{Challenges}

Conceptualizing LLM-as-judge as a measurement instrument (e.g., using an LLM to generate quality ratings for code snippets is akin to measuring the quality of the snippet) gives us a language for evaluating LLM-judges. The LLM-judge should exhibit validity, test-retest reliability, inter-rater reliability, minimal random and systematic error, measurement invariance, etc.~\cite{ralph2024teaching} 

When relying on the judgment of LLMs, researchers must build a \textit{reliable} process for generating judgment labels that considers the non-deterministic nature of LLMs and report the intricacies of that process transparently~\cite{DBLP:journals/corr/abs-2412-12509}.
For example, the order of options has been shown to affect LLM outputs in multiple-choice settings~\cite{DBLP:conf/naacl/PezeshkpourH24}. Additionally, LLMs can behave differently when reviewing their own outputs~\cite{NEURIPS2024_7f1f0218}.

Reliability is a necessary but insufficient condition for validity. 
For example, a reliable LLM might be reliably inaccurate and wrong. Experiments using LLMs to ``peer'' review scholarly articles have not gone well~\cite{DBLP:conf/coling/ZhouC024}. LLM judges are susceptible to adversarial manipulation, where semantics-preserving changes (e.g., misleading code comments) may deceive the judge into accepting flawed artifacts~\cite{DBLP:journals/corr/abs-2510-24367}. For tasks with no single, objectively correct answer, it is not clear whether an LLM judge should reflect the majority opinion or reason  independently. 
The underlying statistical framework of an LLM usually pushes outputs towards the most likely (majority) answer, which might not be the best answer. Questions about bias, accuracy, and trust remain~\cite{DBLP:journals/corr/abs-2406-18403}.

Meanwhile, evaluating and judging large numbers of items, for example, to perform fault localization on the thousands of bugs in large open-source projects, has significant clock time, compute time, and carbon footprint.

Regarding systematic errors, LLMs may suffer from well-documented biases and issues with fairness~\cite{Gallegos2024BiasAF} (e.g., in assessing discussions of pull requests). 

Beyond questions of technical capacity, ethical questions remain, particularly if there is some implicit expectation that a human is judging the output. Involving a human in the judgment loop, for example, to contextualize the scoring, is one approach~\cite{panHumanCenteredDesignRecommendations2024}. 
However, the lack of large-scale ground truth datasets for benchmarking LLM performance in judgment studies hinders progress in this area. There is simply no evidence that LLMs can accurately and reliably judge most relevant properties of SE artifacts, or that offering LLM-based support to human judges improves any specific outcomes.  